We prove necessary residual-dimension conditions for null-projected supersymmetric Yang--Mills constructions with a Majorana--Weyl gaugino under five explicit assumptions on the projection, transformations, translation bracket and absorption of cubic fermion terms. A real gaugino bounds the null-frame rank by the Witt index, and the residual dimension must be two or ten. Consequently the difference between the numbers of spacelike and timelike directions must be zero or eight in absolute value. The main algebraic step is a converse for a cubic gamma-matrix identity, relating the spinor dimension to the vector dimension. This obstructs transferring Nishino's mechanism to the specified higher-dimensional signature families, including nineteen plus three and twenty-seven plus three. These are necessary conditions within the stated axioms; existence in every remaining signature is a separate question. For unitary single-time theories with vanishing tensor charges on a massless multiplet, we also give a helicity-span argument excluding any finite collection of scalar, spinor and vector polarizations when there are more than sixteen real supercharges. All general arguments are written in the text; explicit gamma-matrix computations serve as supplementary checks.
